\subsubsection{Fourth-Order Collision Moment and \texorpdfstring{$\mathrm{R\acute{e}nyi}$}{Renyi}--4 Projection Entropy Rate}\label{subsec:pom-s4}
\begin{definition}[Fourth-order folding collision moment]\label{def:pom-s4}
Define the fourth-order collision sum
$$
S_4(m):=\sum_{x\in X_m} d_m(x)^4.
$$
Equivalently, drawing four windows \(\omega_1,\omega_2,\omega_3,\omega_4\) independently and uniformly from \(\Omega_m\), one has
$$
\mathbb{P}\bigl[\Fold_m(\omega_1)=\Fold_m(\omega_2)=\Fold_m(\omega_3)=\Fold_m(\omega_4)\bigr]=\frac{S_4(m)}{16^m}.
$$
\end{definition}

\begin{proposition}[Five-state nonnegative integer realization of the fourth-order collision kernel]\label{prop:pom-s4-recurrence}
There exists a five-state nonnegative integer matrix
$$
A_4:=
\begin{pmatrix}
0&1&0&0&0\\
0&0&7&0&1\\
0&1&2&0&0\\
1&0&1&0&0\\
0&0&0&1&0
\end{pmatrix},
$$
together with readout vectors
$$
u_4:=(9,404,1613,25,114)^{\mathsf T},\qquad
v_4:=(1,0,0,0,0)^{\mathsf T},
$$
such that for all \(m\ge 2\),
$$
S_4(m)=2\,u_4^{\mathsf T}A_4^{\,m-2}v_4.
$$
Consequently, \(S_4(m)\) satisfies an exact fifth-order recurrence: for all \(m\ge 13\),
$$
S_4(m)=2S_4(m-1)+7S_4(m-2)+2S_4(m-4)-2S_4(m-5).
$$
\end{proposition}
\begin{proof}
By Table \ref{tab:fold_collision_moment_recursions}, one can verify in an auditable manner that \(S_4(m)\) satisfies an integer recurrence of order \(5\) (script \texttt{scripts/exp\_fold\_collision\_moment\_recursions.py}). On the other hand, the above \(A_4\) and \((u_4,v_4)\) provide a completely explicit five-state nonnegative integer realization; its consistency with the recurrence and the enumerated values for small \(m\) can be verified entry by entry via the script \texttt{scripts/exp\_fold\_collision\_kernel\_A4\_verify.py}. \qed
\end{proof}

\begin{corollary}[Fourth-order collision exponent and \texorpdfstring{$h_4$}{h4}]\label{cor:pom-s4-asymptotic}
Let \(r_4:=\rho(A_4)\) be the largest real root of the polynomial
$$
x^5-2x^4-7x^3-2x+2=0
$$
(\(r_4\approx 3.846059041178\)). Then
$$
S_4(m)=\Theta(r_4^m),
\qquad
h_4:=\lim_{m\to\infty}\frac{1}{m}\Bigl(-\log\sum_{x\in X_m}\pi_m(x)^4\Bigr)
:=\log\!\left(\frac{16}{r_4}\right).
$$
\end{corollary}
\begin{proof}
By Proposition \ref{prop:pom-s4-recurrence}, the generating function denominator is \(\det(I-zA_4)\), whose smallest-modulus zero is \(z_\star=1/\rho(A_4)\), whence \(S_4(m)=\Theta(r_4^m)\). Moreover,
$$
\sum_{x\in X_m}\pi_m(x)^4=\sum_x\frac{d_m(x)^4}{16^m}=\frac{S_4(m)}{16^m},
$$
and substituting the asymptotic yields \(h_4=\log(16/r_4)\). \qed
\end{proof}


\begin{remark}[Near-critical RH violation of the fourth-order collision kernel: spectral decomposition numerical certificate (\texorpdfstring{$A_4$}{A4})]\label{rem:pom-collision-rh-break-a4-spectrum}
Write \(r_4=\rho(A_4)\approx 3.846059041178\).
Computing the spectrum directly (equivalently, finding the roots of \(\det(\lambda I-A_4)=\lambda^5-2\lambda^4-7\lambda^3-2\lambda+2\)) yields the non-Perron eigenvalues approximately as
\[
\mu_2\approx -1.963592213449,\qquad
\mu_3\approx 0.503913309522,\qquad
\mu_{4,5}\approx -0.193190068626\pm 0.698726683936\,\mathrm{i}.
\]
Therefore the sub-spectral radius is
\[
\Lambda_4:=\max_{\mu\in\sigma(A_4)\setminus\{r_4\}}|\mu|
:=|\mu_2|
\approx 1.963592213449,
\]
while
\(
\sqrt{r_4}\approx 1.961137180612
\), so the square-root bound of the kernel-version RH (the analogous condition in Definition \ref{def:kernel-rh}) is strictly violated on the fourth-order collision kernel:
\[
\Lambda_4>\sqrt{r_4},\qquad
R_4:=\frac{\Lambda_4}{\sqrt{r_4}}\approx 1.001251841463.
\]
In terms of the ``pole real-part exponent'', this can be written as
\[
\beta_4:=\frac{\log \Lambda_4}{\log r_4}\approx 0.500928740205=\frac12+9.2874\times 10^{-4},
\]
and one may introduce the near-critical offset
\[
\delta_4:=\log\!\left(\frac{\Lambda_4^2}{r_4}\right)=2\log R_4\approx 0.002502117125.
\]
 This excess is dominated by a stable negative real mode (\(\mu_2<0\)), corresponding to an irremovable anti-phase oscillation slow mode.
\end{remark}

\begin{remark}[Single-parameter RH critical point of the fourth-order collision kernel: edge-weight threshold \texorpdfstring{$t_\star$}{t*} (resultant elimination, auditable)]\label{rem:pom-collision-rh-break-a4-threshold}
Treating the unique multi-edge weight \(7\) in \(A_4\) as a continuous parameter \(t>0\), define the one-parameter family
\[
A_4(t):=
\begin{pmatrix}
0&1&0&0&0\\
0&0&t&0&1\\
0&1&2&0&0\\
1&0&1&0&0\\
0&0&0&1&0
\end{pmatrix},
\qquad
p_t(x):=\det(xI-A_4(t))=x^5-2x^4-tx^3-2x+2.
\]
Write \(r(t)=\rho(A_4(t))\), \(\Lambda(t)=\max_{\mu\in\mathrm{spec}(A_4(t))\setminus\{r(t)\}}|\mu|\), and let the Newman-type offset be
\(
\delta(t):=\log(\Lambda(t)^2/r(t))
\)
(consistent with the convention in Remark \ref{rem:pom-collision-rh-break-a4}).
On this family, the critical point where \(\delta(t)\) attains its boundary can be characterized via the \textbf{negative real root channel}: the critical point is equivalent to the existence of \(y<0\) such that \(p_t(y)=p_t(y^2)=0\), whence (excluding the degenerate roots \(y=\pm 1\))
\[
y^8-2y^6+2y^5-2y^4+2y^3-2=0.
\]
Setting \(z:=-y>0\), one equivalently obtains
\(
z^8-2z^6-2z^5-2z^4-2z^3-2=0
\), and from \(p_t(y)=0\) one can explicitly solve for
\[
t=\frac{y^5-2y^4-2y+2}{y^3}
:=z^2+2z-\frac{2}{z^2}-\frac{2}{z^3}.
\]
Further resultant elimination in \(y\) yields an eighth-degree integer-coefficient certificate \(P_4(t)=0\); its unique positive real root is denoted \(t_\star\).
Numerical verification shows: when \(t<t_\star\) one has \(\Lambda(t)<\sqrt{r(t)}\) (strict kernel RH); when \(t=t_\star\) the bound is exactly attained (RH-sharp, with \(\mu_\star=-\sqrt{r_\star}\)); and when \(t>t_\star\) the violation \(\Lambda(t)>\sqrt{r(t)}\) occurs.
 Therefore the integer kernel \(A_4=A_4(7)\) lies on the supercritical side with only a near-critical excess \(7-t_\star\ll 1\), which explains the \(10^{-3}\) magnitude of \(R_4-1\).
\end{remark}
\input{sections/generated/eq_collision_kernel_A4_edge_weight_threshold}


\begin{remark}[Near-RH violation of the fourth-order collision kernel (\texorpdfstring{$A_4$}{A4})]\label{rem:pom-collision-rh-break-a4}
Write \(r_4=\rho(A_4)\), and let the sub-spectral radius be
\[
\Lambda_4:=\max_{\mu\in\mathrm{spec}(A_4)\setminus\{r_4\}}|\mu|.
\]
 Then the fourth-order collision kernel exhibits a \textbf{mild} kernel-version RH excess: \(\Lambda_4>\sqrt{r_4}\) (see Equation \ref{eq:collision_kernel_A4_rh_breaking}), with the corresponding \(s\)-plane nontrivial pole real-part exponent
\(
\beta_4=\log\Lambda_4/\log r_4
\)
 strictly greater than \(1/2\), but only by an excess of order \(10^{-3}\) near criticality.
\par
To uniformly characterize this ``Newman-type offset'', one may define for any finite kernel \(A\)
\[
\delta(A):=\log\!\left(\frac{\Lambda(A)^2}{\rho(A)}\right),
\qquad
\Lambda(A):=\max_{\mu\in\mathrm{spec}(A)\setminus\{\rho(A)\}}|\mu|.
\]
Then the kernel-version RH is equivalent to \(\delta(A)\le 0\). In the present case \(\delta(A_4)>0\) but is numerically very small (Equation \ref{eq:collision_kernel_A4_rh_breaking}).
\par
Combined with Remark \ref{rem:pom-collision-rh-margin} and Remark \ref{rem:pom-collision-rh-margin-a3} (strict margins for \(q=2,3\)), one sees that \(q=4\) is the smallest order of violation within the constructed range.
More generally, for the verified higher-order recurrences (Table \ref{tab:fold_collision_moment_recursions}), a direct root scan of the characteristic polynomials yields the rapid deterioration of
\(\Lambda_q/\sqrt{r_q}\) and \(\beta_q\) as \(q\) increases: the critical line is crossed starting from \(q=4\), and the excess strengthens rapidly for \(q\ge 5\) (Table \ref{tab:fold_collision_kernel_rh_scan_q2_8}).
Furthermore, in the scanned range \(q\le 8\), the largest non-Perron root achieving \(\Lambda_q\) is numerically always a negative real root (anti-phase mode), so the violation is first triggered by a real-axis mode.
\par
\textbf{Testable prediction.} Since \(q=4\) already lies at the near-RH boundary and the dominant sub-mode is a negative real root, one expects that in higher-order \(q\) moment-kernels/recurrences:
(i) the trailing block of recurrence coefficients more closely approaches an alternating boundary pattern (e.g., saturation patterns of \(\pm 2^k\) magnitude);
(ii) the prime factor exponents in Hankel certificates (especially 2-adic and 5-adic) exhibit higher valuations at earlier stages.
This prediction can be directly verified through the recurrence coefficients and Hankel certificate scans for \(A_5\) and beyond.
\end{remark}

\input{sections/generated/eq_collision_kernel_A4_rh_breaking}
\input{sections/generated/tab_fold_collision_kernel_rh_scan_q2_8}


\begin{definition}[\texorpdfstring{$\zeta$}{zeta} function and finite-part constant of the fourth-order collision kernel]\label{def:pom-collision-zeta-a4}
For the five-state kernel \(A_4\), define
$$
\zeta_{A_4}(z):=\exp\!\left(\sum_{n\ge 1}\frac{\mathrm{Tr}(A_4^n)}{n}z^n\right)
::=\frac{1}{\det(I-zA_4)}.
$$
The corresponding primitive orbit counts \(p_n(A_4)\), as well as the Abel/Mertens finite-part constants \(\log\mathfrak{M}(A_4)\) and \(\mathsf{M}(A_4)=\gamma+\log\mathfrak{M}(A_4)\), can be extracted via the universal interface of Section \ref{sec:zeta-finite-part}; the reproducible outputs are given in Table \ref{tab:collision_kernel_A4_primitive} and Table \ref{tab:collision_kernel_A4_finite_part} (scripts \texttt{scripts/exp\_collision\_kernel\_A4\_primitive.py} and \texttt{scripts/exp\_collision\_kernel\_A4\_finite\_part.py}).
\end{definition}

\input{sections/generated/tab_collision_kernel_A4_primitive}
\input{sections/generated/tab_collision_kernel_A4_finite_part}

\par
\medskip\noindent\textbf{Audit conclusion (derivative formula and algebraic certificate for the leading-pole residue constant \(C\): \(A_2,A_3,A_4\))}
Let \(A\) be a finite-state collision kernel (SFT adjacency matrix), with characteristic polynomial normalized to have leading coefficient \(1\):
\[
\chi_A(\lambda):=\det(\lambda I-A),\qquad d:=\deg\chi_A,
\]
and suppose the Perron root \(r=\rho(A)\) is a simple root. Write
\(
\zeta_A(z)=\det(I-zA)^{-1}
\), \(z_\star=1/r\), and following the convention of Tables \ref{tab:collision_kernel_A2_finite_part}--\ref{tab:collision_kernel_A4_finite_part}, define the leading-pole residue constant
\[
C(A):=\lim_{z\to z_\star}(1-rz)\,\zeta_A(z).
\]
Then this constant depends only on \((\chi_A,r)\) and satisfies the closed-form identity
\[
\boxed{\;
C(A)=\frac{r^{d-1}}{\chi_A'(r)}
\;=\;
\prod_{\lambda_j\in\sigma(A)\setminus\{r\}}\left(1-\frac{\lambda_j}{r}\right)^{-1}\; }.
\]
Therefore, once \(\chi_A\) and \(r\) have been independently audited (e.g., given by the recurrence denominator/characteristic polynomial and the Perron root certificate), \(C(A)\) can be exactly reconstructed without computing a residue limit or explicitly finding the full spectrum; moreover, any nonnegative integer realization preserving the same \(\chi_A\) (see Remark \ref{rem:pom-a3-alternative-realization}) must yield the same \(C(A)\).
\par
Furthermore, for the three collision kernels in this paper whose characteristic polynomials have been given:
\[
\chi_{A_2}(\lambda)=\lambda^3-2\lambda^2-2\lambda+2,\qquad
\chi_{A_3}(\lambda)=\lambda^3-2\lambda^2-4\lambda+2,\qquad
\chi_{A_4}(\lambda)=\lambda^5-2\lambda^4-7\lambda^3-2\lambda+2,
\]
the above formula together with elimination allows one to upgrade the three decimal constants \(C(A_q)\) in Tables \ref{tab:collision_kernel_A2_finite_part}--\ref{tab:collision_kernel_A4_finite_part} to integer-coefficient equation-level audit certificates: they respectively satisfy the following minimal polynomials (each having a unique real root in \((0,1)\)):
\[
\boxed{\;
37C^3-37C^2+C+1=0\quad(C=C(A_2)),\qquad
141C^3-141C^2+23C+1=0\quad(C=C(A_3)),}
\]
\[
\boxed{\;
3940876C^5-3940876C^4+920083C^3-33522C^2+168C+4=0\quad(C=C(A_4)).}
\]
This yields a strong audit protocol requiring no decimal alignment: for each \(A_q\), one only needs to publish (i) \(\chi_{A_q}\in\ZZ[\lambda]\), (ii) the minimal polynomial \(P_{C_q}\in\ZZ[C]\) of \(C(A_q)\); together with the \texttt{tail proxy} from the tables as a numerical certificate for the tail truncation error of \(\log\mathfrak{M}\), this upgrades ``reproducibility'' to integer-coefficient equation-level constraints.

\begin{remark}[Constant-level rarity scale: near-halving of \texorpdfstring{$\mathsf{M}(A_4)/\mathsf{M}(A_3)$}{M(A4)/M(A3)}]\label{rem:pom-collision-mertens-halving}
Tables \ref{tab:collision_kernel_A3_finite_part} and \ref{tab:collision_kernel_A4_finite_part} give
\(
\mathsf{M}(A_3)\approx 0.1560083
\)
and
\(
\mathsf{M}(A_4)\approx 0.0730980
\).
Therefore the constant-level contraction ratio is
\[
\boxed{\;
\frac{\mathsf{M}(A_4)}{\mathsf{M}(A_3)}\approx 0.4686\;}
\]
 which is numerically close to one-half.
This ratio provides a \textbf{fine-structure fingerprint independent of the exponential main scale \(r_q\)}: when the collision multiplicity is raised from triple to quadruple,
besides the change in the exponential scale given by the Perron root \(r_q\), the constant-level drift of the primitive leading term (Mertens constant) also contracts significantly.
Moreover, the extremely small \texttt{tail proxy} in the tables indicates that the extraction of this finite-part constant is numerically highly stable, making it suitable as a cross-order rarity scale.
\end{remark}

\begin{remark}[Step-size acceleration of the finite-part constant (\(A_2\to A_4\))]\label{rem:pom-collision-mertens-accel}
From Tables \ref{tab:collision_kernel_A2_finite_part}--\ref{tab:collision_kernel_A4_finite_part},
\(
\mathsf{M}(A_2)\approx 0.2325880,\,
\mathsf{M}(A_3)\approx 0.1560083,\,
\mathsf{M}(A_4)\approx 0.0730980
\).
Therefore
\[
\Delta_3:=\mathsf{M}(A_2)-\mathsf{M}(A_3)\approx 0.07658,\qquad
\Delta_4:=\mathsf{M}(A_3)-\mathsf{M}(A_4)\approx 0.08291.
\]
One sees that at \(q=2,3,4\) the decay step size of the finite-part constant exhibits an increasing trend, providing a fine-structure fingerprint that appears in parallel with the near-RH critical line behavior.
\end{remark}
